% Emoji 和特殊字符字体支持
\usepackage{newunicodechar}
\newfontfamily\DejaSans{Segoe UI Emoji}
\newfontfamily\TreeFont{DejaVu Sans Mono}
\newfontfamily\CircleFont{Microsoft YaHei}

% 定义 emoji 字符
\newunicodechar{💻}{{\DejaSans 💻}}
\newunicodechar{✅}{{\DejaSans ✅}}
\newunicodechar{🌟}{{\DejaSans 🌟}}
\newunicodechar{📘}{{\DejaSans 📘}}
\newunicodechar{🔗}{{\DejaSans 🔗}}
\newunicodechar{📝}{{\DejaSans 📝}}
\newunicodechar{💪}{{\DejaSans 💪}}
\newunicodechar{📚}{{\DejaSans 📚}}
\newunicodechar{💡}{{\DejaSans 💡}}
\newunicodechar{📖}{{\DejaSans 📖}}
\newunicodechar{⭐}{{\DejaSans ⭐}}
\newunicodechar{🤖}{{\DejaSans 🤖}}
\newunicodechar{🎯}{{\DejaSans 🎯}}
\newunicodechar{🚀}{{\DejaSans 🚀}}
\newunicodechar{🎉}{{\DejaSans 🎉}}
\newunicodechar{🟩}{{\DejaSans 🟩}}
\newunicodechar{🟦}{{\DejaSans 🟦}}
\newunicodechar{📷}{{\DejaSans 📷}}
\newunicodechar{✔}{{\DejaSans ✔}}
\newunicodechar{✓}{{\DejaSans ✓}}
\newunicodechar{✗}{{\DejaSans ✗}}

% 定义带圈数字 - 使用 CircleFont
\newunicodechar{①}{{\CircleFont ①}}
\newunicodechar{②}{{\CircleFont ②}}
\newunicodechar{③}{{\CircleFont ③}}
\newunicodechar{④}{{\CircleFont ④}}
\newunicodechar{⑤}{{\CircleFont ⑤}}
\newunicodechar{⑥}{{\CircleFont ⑥}}
\newunicodechar{⑦}{{\CircleFont ⑦}}
\newunicodechar{⑧}{{\CircleFont ⑧}}
\newunicodechar{⑨}{{\CircleFont ⑨}}
\newunicodechar{⑩}{{\CircleFont ⑩}}

% 定义树形图字符（目录结构字符）
\newunicodechar{─}{{\TreeFont ─}}
\newunicodechar{│}{{\TreeFont │}}
\newunicodechar{├}{{\TreeFont ├}}
\newunicodechar{└}{{\TreeFont └}}

% 移除章节编号
\setcounter{secnumdepth}{0}

% 设置更大的字体
\usepackage{geometry}
\geometry{a4paper, margin=1in}

% 设置正文字体大小为 12pt
\renewcommand{\normalsize}{\fontsize{12pt}{18pt}\selectfont}

% 超链接样式 - 蓝色高亮带下划线
\usepackage{hyperref}
\hypersetup{
    colorlinks=true,
    linkcolor=blue,
    filecolor=blue,
    urlcolor=blue,
    citecolor=blue,
    anchorcolor=blue,
    bookmarks=true,
    pdfborder={0 0 0}
}

% 为超链接添加下划线
\usepackage[normalem]{ulem}
\makeatletter
\Hy@colorlinkstrue
\def\Hy@colorlink#1{%
  \begingroup
    \color{#1}%
    \uline
}
\def\Hy@endcolorlink{%
  \endgroup
}
\makeatother

% 代码块样式 - 带边框和语法高亮
\usepackage{listings}
\usepackage{xcolor}

% 定义颜色
\definecolor{codebackground}{RGB}{248,248,248}
\definecolor{codeborder}{RGB}{204,204,204}
\definecolor{keywordcolor}{RGB}{0,0,255}
\definecolor{commentcolor}{RGB}{0,128,0}
\definecolor{stringcolor}{RGB}{163,21,21}
\definecolor{numbercolor}{RGB}{128,128,128}

% 配置 listings 全局样式
\lstdefinestyle{mystyle}{
    backgroundcolor=\color{codebackground},
    commentstyle=\color{commentcolor}\itshape,
    keywordstyle=\color{keywordcolor}\bfseries,
    numberstyle=\tiny\color{numbercolor},
    stringstyle=\color{stringcolor},
    basicstyle=\ttfamily\normalsize,
    breakatwhitespace=false,
    breaklines=true,
    captionpos=b,
    keepspaces=true,
    showspaces=false,
    showstringspaces=false,
    showtabs=false,
    tabsize=4,
    frame=single,
    frameround=tttt,
    rulecolor=\color{codeborder},
    framesep=4pt,
    xleftmargin=8pt,
    xrightmargin=8pt,
    aboveskip=10pt,
    belowskip=10pt
}

\lstset{style=mystyle}

% 为不同语言配置高亮
\lstdefinelanguage{makefile}{
    morekeywords={CC,CFLAGS,all,clean},
    morecomment=[l]{\#},
    morestring=[b]",
}

% 设置默认语言
\lstset{language=C}
